\documentclass{article}
\usepackage{amsmath,amssymb}
\begin{document}
\section*{Physical positive controls for the four-copy AC contraction equations}
These are numerical consistency checks, separate from the exact rational
nonexistence certificate. They verify that the same physical contraction
model admits known fourth-order AC states at neighboring sizes.

On the normalized occupation basis of $\operatorname{Sym}^4(V_{N/2})$,
annihilation of component $i$ is normalized as
\[
 A_i|\nu\rangle=\sqrt{\nu_i/4}\,|\nu-e_i\rangle.
\]
For a pure tensor $|\psi\rangle^{\otimes4}$ this gives
$A_i|\psi\rangle^{\otimes4}=c_i|\psi\rangle^{\otimes3}$.
Contracting against a rank-$L$, component-zero spin tensor gives
\[
 \sum_i(T_{L0})_{ii}\,A_i X A_i^\dagger=0,
 \qquad L=1,2,3,4,
\]
when $X$ is the rotational average of a fourth-order AC pure tensor.
The real part of this rotational average is positive semidefinite and
still satisfies these real equations.

The verification explicitly decomposes the four-copy physical space into
spin representations, forms the rotational average of each known state,
and tests its trace and every contraction equation. To prevent roundoff
leakage during long lowering ladders, each descent is projected onto the
known total-spin Casimir eigenspace and polar-reorthonormalized. Numerical
matrix rank is never used as the nonexistence proof.

The input at $N=12$ is
\[
\frac{-\sqrt7|D_{12}^{(1)}\rangle+\sqrt{11}|D_{12}^{(6)}\rangle
+\sqrt7|D_{12}^{(11)}\rangle}{5}.
\]
At $N=14$ the input is the exact radical witness saved in
\texttt{../n14-t4/basis.json}. Both have directly verified
$\rho_4=I_5/5$.
The numerical invariant systems have respectively $250$ and $447$
real symmetric block variables, ranks $242$ and $417$, and nullities
$8$ and $30$. The explicit state twirls have trace one and contraction
residuals below $2\cdot10^{-16}$ in Euclidean norm.

For $N=13$, the stabilized physical system has $331$ variables,
rank $320$, and nullity $11$. Its trace functional vanishes on that
nullspace numerically. The decisive statement is the separate exact
rational verification of this trace annihilation, not this numerical rank
or a failed optimization. The positive controls show that the construction
does not exclude the known adjacent examples.
\end{document}
